% TODO: AI - Write abstract (150-250 words)
% Key points to include:
% - Background on respiratory pathogen seasonality
% - Gap in understanding concurrent outbreaks
% - Methodology: survival analysis approach
% - Key findings: geographic differences, threshold effects
% - Clinical/public health implications

The seasonal dynamics of respiratory pathogens such as influenza and respiratory syncytial virus (RSV) exhibit complex temporal patterns that vary geographically. While individual pathogen seasonality has been extensively studied, the understanding of concurrent multi-pathogen outbreaks remains limited. This study employs a novel survival analysis approach to characterize the timing and frequency of joint respiratory disease outbreaks across different geographic regions and risk thresholds. Using simulation data from 1000 model runs across eight Chinese cities, we analyzed the cumulative incidence of concurrent outbreaks at multiple severity thresholds. Our findings reveal significant geographic variation in outbreak timing, with northern cities showing earlier and more frequent joint outbreaks compared to southern regions. The survival analysis approach properly handles right-censoring of outbreak-free periods and provides robust estimates of outbreak intervals. These results have important implications for public health preparedness and resource allocation during respiratory disease seasons.



